% DRAFT SKELETON — Phase 3 prep, numbers pending federal run chain (Phase 1/2). NOT \input anywhere.
% ============================================================================
% Standalone draft of the proposed main-text section "The Audit on a Second
% Platform". Placement is decision D-iii (pending). DO NOT \input. DO NOT edit
% values.tex; every numeric claim below is a \valTODO placeholder carrying an
% inline % TODO_NUMBER tag for the lead to wire into values.tex after the
% federal run chain completes. Voice, macro style, and table conventions mirror
% sec04_validation_submission.tex and sec05_cobidder_type_submission.tex.
%
% Placeholder macro convention used here (defined in DRAFT_MACROS block below
% only so this file compiles standalone; the lead should REPLACE these with
% real values.tex bindings and DELETE the block before integration):
% ----------------------------------------------------------------------------
\providecommand{\valTODO}{\textbf{[??]}}
% ============================================================================

\section{The Audit on a Second Platform}
\label{sec:comparative}

The validation in Section~\ref{sec:validation} draws a boundary on
one platform: under a reproducible, non-circular label, the
award-layer score concentrates adjudication-anchored loser-side
exposure in the raw data, but almost nothing survives once
procurement opportunity is held fixed, and the prospective ranking
collapses outside the incumbent pool. The defensible product is not
the screen but the audit---the decomposition that tells an agency
how far a cheap award-layer statistic can be trusted before costly
bid-layer work begins. A decomposition is only worth carrying,
however, if it travels. This section runs the same audit, unchanged,
on a second procurement system, and asks what the protocol returns
where the institutions differ. The exercise tests portability of the
\emph{method} and of the loser-side construct; it is not an
independent cartel universe, and we read it as neither confirmation
nor refutation of the legal status of any firm.

\subsection{A Different Procurement System}
\label{sec:comparative_setting}

The federal platform differs from BEC--SP on the dimensions that
matter for the audit. ComprasNet is a single national system: where
BEC is one state-level platform inside the federal legal frame,
ComprasNet aggregates procuring units (UASGs) across the whole
federal administration. Its modality mix is also narrower. The
federal panel is effectively pure Preg\~ao---about $15\%$ regular
electronic auction and $85\%$ price-registration (SRP, \emph{sistema
de registro de pre\c{c}os})---and the sealed-bid Convite modality
that supplies the institutional contrast in the main text is
extinct federally. The audit therefore runs on a panel with no
modality margin to stratify on, which sharpens rather than weakens
the test: the loser-side construct must earn its keep without the
modality variation that organized the BEC reading. The observable
window is $2013$--$2019$, the federal-overlap window with the BEC
panel.

The observability gap is the second contrast, and it is on-thesis.
The public federal data expose participation and the winner flag,
but no bid microdata---there is no federal analogue to the BEC
\texttt{LANCES} bid ladder. A screen that reads only the award layer
is exactly the screen one \emph{can} build here; the bid-layer
forensic step that the main text positions downstream of the screen
cannot be constructed from the public federal release at all. The
federal panel is thus the natural stress test for a cheap award-layer
audit: it is the environment where bid-microdata screens are
unavailable, so the question of how far the award layer alone can be
trusted is not academic but binding.

\subsection{Data, Label, and Linkage}
\label{sec:comparative_data}

The federal panel comprises $\valTODO$ participation rows % TODO_NUMBER \valFedPanelRows (51.0M)
across $\valTODO$ distinct CNPJ-14 firms, % TODO_NUMBER \valFedFirms (92,600)
built from the Comptroller-General (CGU) open bulk dumps and
restricted to the Preg\~ao modalities (regular and SRP) under the
canonical schema used throughout the main text. Of these,
$\valTODO$ are always-losers---firms with a win rate of zero over % TODO_NUMBER \valFedAlwaysLosers (35,943)
the observed window. Applying the \emph{same} IQR rule used for
BEC---$\textit{median} + 1.5\cdot\textit{IQR}$ on the tenders-count
distribution of always-losers, not the Tukey upper-quartile
rule---re-derives a federal frequent-loser threshold of $\valTODO$ % TODO_NUMBER \valFedFLthreshold (32)
(against $14$ on BEC), and the canonical federal frequent-loser
count above that threshold is $\valTODO$ firms. The threshold is % TODO_NUMBER \valFedFLcount (6,491)
re-estimated from the federal data, not transported: the construct
is the rule, not the cutoff value, exactly as the rules of
engagement in the main text require.

The legal anchors are the \emph{same} CADE cases that anchor the BEC
validation---$\valTODO$ numbered adjudicated procurement-cartel % TODO_NUMBER \valFedAudCases (7)
processes, with one unnumbered summary case excluded for lack of
firm-level identification. This is the honest statement of what the
federal exercise is: a second platform with \emph{partially
overlapping} legal anchors, not a fresh cartel population. The
federal evidence therefore tests whether the audit protocol and the
loser-side construct \emph{port} across procurement systems; it does
not assemble an independent ground truth, and we do not read it as
one. Direct defendants are matched to federal participant records by
the canonical establishment-anchored CNPJ rule. Expanding the
anchored tender-items yields $\valTODO$ federal cobidders under the % TODO_NUMBER \valFedAudCobiddersAll (3,851)
broad rule and $\valTODO$ under the broad always-loser restriction % TODO_NUMBER \valFedAudCobiddersBroadAL (196)
that mirrors the main-text validation target; the frequent-loser
flag is never used to build the label, so the federal validation is
non-circular in the score by the same construction as the BEC
validation.

\subsection{Side-by-Side Audit Verdict}
\label{sec:comparative_verdict}

Table~\ref{tab:comparative_audit} runs each stage of the
Section~\ref{sec:validation} battery on both platforms and reports
the comparable quantity in each column. The point of the table is
not the absolute levels but the \emph{shape} of the audit: whether
each successive discipline---opportunity adjustment, within-stratum
restriction, power-bounding, matched permutation, label-frozen
timing, case concentration, negative controls---moves the federal
verdict the way it moves the BEC verdict.

\begin{table}[htbp]
\centering
\caption{The Audit on Two Platforms: BEC--SP versus ComprasNet}
\label{tab:comparative_audit}
\scriptsize
\begin{adjustbox}{max width=\textwidth,center}
\begin{threeparttable}
\begin{tabular}{p{5.4cm}cc}
\toprule
Audit stage & BEC--SP & ComprasNet \\
\midrule
Raw award-layer ROC-AUC & $\valTODO$ & $\valTODO$ \\ % TODO_NUMBER \valExpRawScoreAUC | \valFedAudRawAUC
Label-blind opportunity expectation (AUC) & $\valTODO$ & $\valTODO$ \\ % TODO_NUMBER \valArmorExpLB | \valFedAudLabelBlindAUC
Within-stratum residual (AUC, MEDIUM cells) & $\valTODO$ & $\valTODO$ \\ % TODO_NUMBER \valExpWithinAUC | \valFedAudWithinAUC
Power-bounded residual (det.\ prob.\ @ true AUC $0.55$) & $\valTODO$ & $\valTODO$ \\ % TODO_NUMBER \valArmorPowerFive | \valFedAudPowerFive
Matched-permutation $p$ (within-strata shuffle) & $\valTODO$ & $\valTODO$ \\ % TODO_NUMBER \valExpPermCp | \valFedAudPermCp
Label-frozen prospective timing (AUC) & $\valTODO$ & $\valTODO$ \\ % TODO_NUMBER \valArmorFrozenProspAUC | \valFedAudFrozenProspAUC
Top-case concentration (share of positives) & $\valTODO$ & $\valTODO$ \\ % TODO_NUMBER \valTopCasePosShare | \valFedAudTopCaseShare
Negative-control verdict & $\valTODO$ & $\valTODO$ \\ % TODO_NUMBER \valTODO(BEC: generic geometry) | \valFedAudNegControlVerdict
\bottomrule
\end{tabular}
\begin{tablenotes}
\scriptsize
\item \textit{Notes:} Each row is the comparable quantity from the
Section~\ref{sec:validation} battery, re-run on the federal panel
under the same definitions; the federal construction is detailed in
\ref{app:federal_audit_submission}. Cobidder labels are
adjudication-anchored exposure, not cartel membership, on both
platforms; neither column ranks defendants. The federal panel has no
sealed-bid modality, so no modality stratification is reported.
PR-AUC is the lead metric for the rare-target rows; ROC-AUC is shown
for comparability. Every cell is a placeholder pending the federal
run chain (\valTODO).
\end{tablenotes}
\end{threeparttable}
\end{adjustbox}
\end{table}

\subsection{Interpretation}
\label{sec:comparative_interpretation}

% ----------------------------------------------------------------------------
% CONDITIONAL DRAFT — two mutually exclusive paragraphs. After the federal run
% chain resolves the within-stratum / permutation rows of
% Table~\ref{tab:comparative_audit}, KEEP exactly one and DELETE the other.
% Both obey the locked language rules: flags / screens / concentrates risk;
% never detects / proves / outperforms. Neither makes a legal claim.
% ----------------------------------------------------------------------------

% --- DRAFT A: deflation replicates (federal residual also vanishes) ----------
\paragraph{[Draft A --- deflation replicates.]}
On the federal panel the audit returns the same verdict it returns
on BEC. The raw award-layer score concentrates loser-side exposure
($\valTODO$ ROC-AUC), but the concentration is again what % TODO_NUMBER \valFedAudRawAUC
procurement opportunity already implies: the label-blind opportunity
ranking reaches $\valTODO$ on its own, the within-stratum residual % TODO_NUMBER \valFedAudLabelBlindAUC
collapses to $\valTODO$, and the matched-strata permutation does not % TODO_NUMBER \valFedAudWithinAUC
reject ($p=\valTODO$). The deflation is platform-agnostic: the cheap % TODO_NUMBER \valFedAudPermCp
screen's apparent reach is opportunity arithmetic on the federal
system as well. What ports is the audit, not a deployable ranking.
An agency on either platform that runs this decomposition learns,
before opening any bid record, that loser-side concentration should
not be used as a stand-alone prioritization device here---and on the
federal platform that lesson is sharper still, because the bid-layer
step that might follow the screen cannot be constructed from public
data at all.

% --- DRAFT B: residual survives federally (institution-coupled value) --------
\paragraph{[Draft B --- residual survives.]}
On the federal panel the audit returns a \emph{different} verdict at
the margin that matters. The raw concentration is again largely
opportunity arithmetic---the label-blind opportunity ranking reaches
$\valTODO$ on its own---but unlike BEC a residual ordering survives % TODO_NUMBER \valFedAudLabelBlindAUC
the within-stratum restriction ($\valTODO$, above the BEC chance % TODO_NUMBER \valFedAudWithinAUC
level), the matched-strata permutation rejects ($p=\valTODO$), and % TODO_NUMBER \valFedAudPermCp
the power calibration places this above the threshold the test can
detect. We read this as institution-coupled screening value, not as
detection: where the federal system removes the modality and
bid-layer information that a richer platform exposes, a cheap
loser-side screen retains residual concentration that the audit on
BEC dissolves. The construct flags risk more informatively exactly
where the cheaper information set binds. The same audit thus does
double duty---it deflates the screen on one platform and locates
residual value on another---which is the property a portable
protocol should have.

% ----------------------------------------------------------------------------

\subsection{Observability as an On-Thesis Finding}
\label{sec:comparative_observability}

The two platforms differ in what they let an analyst see, and the
differences are not nuisance---they are the finding. Item groups are
not observed federally in the BEC sense: federal item codes are
embedded in buyer-specific catalogs rather than a shared
classification, so the item-group margin that the BEC audit
stratifies on (the COARSE cell) has no clean federal counterpart,
and the federal opportunity cells are built on the buyer$\times$year
margins instead (\ref{app:federal_audit_submission}). The bid tier
is absent entirely: no federal bid ladder means no bid-dispersion
benchmark, no bid-layer forensic comparator, and no way to run the
downstream costly-proof step on public data. Both gaps point the
same way. The award-layer audit is precisely the audit that remains
feasible where the richer information is missing, which is the
common case for an enforcement agency staring at a national platform
with participation records and little else. That the loser-side
construct can be re-derived, the label re-anchored, and the full
deflation battery re-run on a system with no bid microdata is itself
the portability result: the cheap screen runs where the expensive
screens cannot be built, and the audit tells the agency how far to
trust it before the expensive screens---if they ever become
available---are worth commissioning.
